\section{Orbitwise arithmetic and a fixed-plane MARKED control}
\label{sec:arithmetic}

The global trace coefficients have arithmetic structure at regular words,
but the character generally depends on the orbit.  This section separates
that theorem from a full-Rauzy fixed-plane stress test.

\subsection{Quadratic prime series at regular words}

Fix a word \(w\) satisfying

\begin{equation}
 D_w=\det(I-g_w)\ne0.
 \label{eq:D-nonzero}
\end{equation}

Let \(\varepsilon_w\) be the quadratic Kronecker character determined by the square
class of \(D_w\).  The regular case of the finite-Weil character formula
gives

\begin{equation}
 \Theta_p(g_w)=\varepsilon_w(p)
 \quad\text{for every odd }p\nmid D_w.
 \label{eq:good-prime-quadratic}
\end{equation}

No quadratic \(L\)-data are assigned here when \(D_w=0\).

For a quadratic Dirichlet character \(\chi\), define, initially for
\(\operatorname{Re}\zeta>1\),

\begin{equation}
 P_\chi(\zeta)=\sum_p\chi(p)p^{-\zeta},
 \qquad
 P_\chi^{\odd}(\zeta)
 =P_\chi(\zeta)-\chi(2)2^{-\zeta}.
 \label{eq:odd-prime-series}
\end{equation}

Then \cref{eq:good-prime-quadratic} gives the exact decomposition

\begin{equation}
 \sum_{p\ \odd}p^{-\zeta}\Theta_p(g_w)
 =P_{\varepsilon_w}^{\odd}(\zeta)+E_w(\zeta),
 \label{eq:word-prime-series}
\end{equation}

where

\begin{equation}
 E_w(\zeta)=
 \sum_{\substack{p\ \odd\\p\mid D_w}}
 \bigl(\Theta_p(g_w)-\varepsilon_w(p)\bigr)p^{-\zeta}
 \label{eq:finite-correction}
\end{equation}

is a finite singular-prime correction.

The Euler logarithm gives an optional expression in familiar \(L\)-data.
For \(\operatorname{Re}\zeta>1\), fix \(\Log L(\zeta,\chi)\) by its
absolutely convergent Euler series at infinity.  M\"obius inversion yields

\begin{equation}
 P_\chi(\zeta)
 =\sum_{k\ge1}\frac{\mu(k)}{k}
   \Log L(k\zeta,\chi^k).
 \label{eq:prime-L-inversion}
\end{equation}

Here \(\chi^k\) means the pointwise power and may be imprimitive.  Formula
\eqref{eq:prime-L-inversion} follows by expanding
\(\Log L(\zeta,\chi)=\sum_{m\ge1}P_{\chi^m}(m\zeta)/m\) and applying the
divisor identity \(\sum_{d\mid n}\mu(d)=0\) for \(n>1\); all rearrangements
are absolutely convergent in the stated half-plane~\citep{apostol1976}.  We
make no logarithm choice across an \(L\)-zero and claim no continuation of
the prime series.

The character \(\varepsilon_w\) is orbit-dependent.  A bounded census in the local
finite-Weil stage found 150 distinct arithmetic signatures among 150 sampled
induced branches.  This is exact finite evidence of fragmentation in that
window, not an all-length theorem.  Repetition may change \(D_w\), and the
correct factor remains \(\Theta_p(g_w^r)\), not \(\varepsilon_w(p)^r\).

\subsection{A general marked convergence threshold}

For this subsection, let \(J_\star\) be any frozen unimodular integral
symplectic form of rank four.  Write \(\rho_{p,\star}\) for the full finite
Weil representation pulled back to that frame and
\(\Theta_{p,\star}(g)=\Tr\rho_{p,\star}(g\bmod p)\).  The proof of
\cref{lem:rank-stability} uses only integral-minor rank stability and the
Thomas magnitude formula, so it applies verbatim in this frozen frame.

\begin{proposition}[Marked absolute threshold]
\label{prop:marked-threshold}
Fix \(g\in\Sp(J_\star,\Z)\) and put

\[
 k_\Q=\dim_\Q\ker(g-I).
\]

For complex \(\alpha\),

\begin{equation}
 \sum_{p\ \odd}|p^{-\alpha}\Theta_{p,\star}(g)|<\infty
 \quad\Longleftrightarrow\quad
 \operatorname{Re}\alpha>1+\frac{k_\Q}{2}.
 \label{eq:marked-threshold}
\end{equation}
\end{proposition}

\begin{proof}
By \cref{lem:rank-stability}, outside finitely many primes,

\[
 |\Theta_{p,\star}(g)|=p^{k_\Q/2}.
\]

Since \(|p^{-\alpha}|=p^{-\operatorname{Re}\alpha}\), the absolute series
is, up to finitely many terms, the prime series with exponent
\(\operatorname{Re}\alpha-k_\Q/2\).  Its convergence criterion and boundary
divergence give \cref{eq:marked-threshold}.
\end{proof}

\subsection{The C24-P073 full-Rauzy MARKED control}

The fixed-plane control comes from the frozen C24 ledger of eventually
positive \emph{full labeled Rauzy cycles}.  It is not an induced C26 branch.
In this paragraph \(\Theta_{p,\mathrm{C24}}\) denotes the finite-Weil
character pulled back to the frozen C24 symplectic frame.
The base-trivialized matrix of cycle P073 is

\begin{equation}
 g_{073}=
 \begin{pmatrix}
 1&8&4&4\\
 1&13&6&6\\
 1&6&4&3\\
 1&2&1&2
 \end{pmatrix},
 \label{eq:P073-matrix}
\end{equation}

and

\begin{equation}
 \det(xI-g_{073})=(x-1)^2(x^2-18x+1).
 \label{eq:P073-charpoly}
\end{equation}

All \(3\times3\) minors of \(g_{073}-I\) vanish, while the gcd of its
\(2\times2\) minors equals one.  Hence \(g_{073}-I\) has rank two over every
finite field.  In the frozen C24 symplectic frame the two-dimensional Thomas
quotient form is represented by

\begin{equation}
 Q_{073}=\begin{pmatrix}1&0\\8&-4\end{pmatrix}.
 \label{eq:P073-Q}
\end{equation}

The exact character formula therefore gives, for every odd prime,

\begin{equation}
 \Theta_{p,\mathrm{C24}}(g_{073})
 =\left(\frac{-4}{p}\right)
  \left(\frac{-1}{p}\right)p=p.
 \label{eq:P073-character}
\end{equation}

Consequently the dimension-normalized \emph{marked} contribution diverges:

\begin{equation}
 \sum_{p\ \odd}\frac{\Theta_{p,\mathrm{C24}}(g_{073})}{p^2}
 =\sum_{p\ \odd}\frac1p=\infty.
 \label{eq:P073-marked-divergence}
\end{equation}

Among the 146 frozen C24 cycles, the exact fixed-space census is \(125\),
\(20\), and \(1\) in dimensions zero, one, and two; P073 is the unique
dimension-two entry.  Thus
\cref{eq:P073-marked-divergence} closes the
\textbf{full-C24 dimension-normalized MARKED assembly}.  It does not show
that the positive-prefix C26 induced language contains a fixed-plane word.
The all-word fixed-plane gate for that narrower language remains open and is
not needed for the sharp Schatten or damped determinant theorems.
